
\documentclass[11pt,a4paper]{article}
\usepackage[utf8]{inputenc}
\usepackage[T1]{fontenc}
\usepackage[french]{babel}
\usepackage{lmodern}
\usepackage{microtype}
\usepackage[a4paper,margin=1.65cm,headheight=15pt]{geometry}
\usepackage{amsmath,amssymb,mathtools}
\usepackage{array,tabularx,multicol,booktabs}
\usepackage{enumitem}
\usepackage[most]{tcolorbox}
\usepackage{xcolor}
\usepackage{tikz}
\usetikzlibrary{arrows.meta,calc,positioning}
\usepackage{pdflscape}
\usepackage{fancyhdr}
\usepackage{hyperref}
\usepackage{pagecolor}
\usepackage{needspace}

\definecolor{fondpage}{RGB}{247,248,250}
\definecolor{encre}{RGB}{31,35,40}
\definecolor{bleu}{RGB}{40,91,135}
\definecolor{vert}{RGB}{55,125,92}
\definecolor{orange}{RGB}{178,105,42}
\definecolor{violet}{RGB}{101,77,145}
\definecolor{rouge}{RGB}{170,72,72}
\definecolor{gris}{RGB}{86,91,99}
\definecolor{bleufonce}{RGB}{28,55,120}
\definecolor{bleuclair}{RGB}{238,243,255}
\definecolor{grisclair}{RGB}{245,245,245}

\hypersetup{colorlinks=true,linkcolor=bleufonce,urlcolor=bleufonce,
pdftitle={Parcours de revision -- Variables aléatoires -- Premiere specialite},pdfauthor={}}

\setlength{\parindent}{0pt}
\setlength{\parskip}{0.25em}
\renewcommand{\arraystretch}{1.13}
\setlength{\tabcolsep}{4pt}
\setlist[itemize]{leftmargin=1.25em,itemsep=0.15em,topsep=0.2em}
\setlist[enumerate,1]{label=\textbf{\arabic*.}, leftmargin=1.05cm, itemsep=0.35em, topsep=0.2em}
\setlist[enumerate,2]{label=\textbf{\alph*.}, leftmargin=0.85cm, itemsep=0.25em, topsep=0.15em}

\pagestyle{fancy}
\fancyhf{}
\cfoot{\thepage}

\newcommand{\R}{\mathbb{R}}
\newcommand{\N}{\mathbb{N}}
\newcommand{\sourceq}[1]{ {\scriptsize\textcolor{bleufonce}{(site Premiere q#1)}}}
\newcounter{question}
\newcounter{reponse}

\newenvironment{blocquestions}[1]{%
  \begin{tcolorbox}[enhanced,breakable,colback=bleuclair,colframe=bleufonce,boxrule=0.6pt,arc=2mm,left=2mm,right=2mm,top=2mm,bottom=1.5mm,title={\bfseries #1},coltitle=white,colbacktitle=bleufonce,before skip=3mm,after skip=3mm, fontupper=\large]
  \large
  \begin{enumerate}[leftmargin=*,label=\textbf{\arabic*.},itemsep=0.82em,topsep=0.5em]
  \setcounter{enumi}{\value{question}}
}{%
  \setcounter{question}{\value{enumi}}
  \end{enumerate}
  \end{tcolorbox}
}

\newenvironment{blocreponses}[1]{%
  \begin{tcolorbox}[enhanced,breakable,colback=bleuclair,colframe=bleufonce,boxrule=0.6pt,arc=2mm,left=2mm,right=2mm,top=2mm,bottom=1.5mm,title={\bfseries #1},coltitle=white,colbacktitle=bleufonce,before skip=3mm,after skip=3mm, fontupper=\large]
  \large
  \begin{enumerate}[leftmargin=*,label=\textbf{\arabic*.},itemsep=0.42em,topsep=0.5em]
  \setcounter{enumi}{\value{reponse}}
}{%
  \setcounter{reponse}{\value{enumi}}
  \end{enumerate}
  \end{tcolorbox}
}

\newtcolorbox{correctionbox}[1]{enhanced,breakable,colback=grisclair,colframe=black!55,boxrule=0.55pt,arc=2mm,left=2mm,right=2mm,top=2mm,bottom=1.5mm,title=\textbf{#1},coltitle=black,colbacktitle=black!12,before skip=2mm,after skip=2mm}
\newtcolorbox{methodebox}[1]{enhanced,breakable,colback=white,colframe=bleu!70!black,boxrule=0.55pt,arc=2mm,left=2mm,right=2mm,top=2mm,bottom=1.5mm,title=\textbf{#1},coltitle=white,colbacktitle=bleu!70!black,before skip=2mm,after skip=2mm}

\newcommand{\titreSujet}[2]{%
  \subsection*{#1}
  \addcontentsline{toc}{subsection}{#1}
  \textit{Référence / inspiration : #2}\par\medskip\hrule\medskip
}
\newcommand{\qcm}[4]{%
\begin{center}\begin{tabularx}{0.96\linewidth}{@{}XX@{}}
\textbf{a.} #1 & \textbf{b.} #2\\[1mm]
\textbf{c.} #3 & \textbf{d.} #4
\end{tabularx}\end{center}}

\tcbset{enhanced,arc=2mm,outer arc=2mm,boxrule=0.42pt,left=1.15mm,right=1.15mm,top=0.75mm,bottom=0.75mm,boxsep=0.35mm,before skip=0.8mm,after skip=0.8mm,fonttitle=\bfseries\footnotesize,coltitle=encre,before upper=\raggedright\fontsize{7.15}{7.65}\selectfont}
\newtcolorbox{fichebox}[3][]{colback=white,colframe=#2!72!black,colbacktitle=#2!18,coltitle=#2!50!black,borderline west={0.9mm}{0pt}{#2!78!black},title={#3},drop fuzzy shadow=black!5,#1}
\newcommand{\tagcol}[2]{\begin{tcolorbox}[enhanced,colback=#1!11,colframe=#1!45!black,boxrule=0.42pt,arc=2mm,left=1mm,right=1mm,top=0.45mm,bottom=0.45mm,before skip=0mm,after skip=0.85mm,fontupper=\bfseries\scriptsize,colupper=#1!55!black]\centering #2\end{tcolorbox}}
\newtcolorbox{panel}[1]{colback=fondpage,colframe=fondpage,boxrule=0pt,arc=0mm,left=0mm,right=0mm,top=0mm,bottom=0mm,height=168mm,valign=top,before skip=0mm,after skip=0mm}

\newcommand{\loitableau}{%
\begin{center}
\begin{tabular}{c|cccc}
$x_i$ & $x_1$ & $x_2$ & $\cdots$ & $x_k$ \\ \hline
$P(X=x_i)$ & $p_1$ & $p_2$ & $\cdots$ & $p_k$
\end{tabular}
\end{center}}

\newcommand{\schemaMoyenne}{%
\begin{center}
\begin{tikzpicture}[scale=0.62,>=Latex]
\draw[->,thick,gris] (-2.4,0)--(2.6,0);
\foreach \x/\lab in {-1.7/$x_1$, -.4/$x_2$, 1.3/$x_3$}{\draw[thick] (\x,.08)--(\x,-.08) node[below] {\scriptsize \lab};}
\draw[very thick,rouge!80!black] (.25,.28)--(.25,-.28) node[above=3pt] {\scriptsize $E(X)$};
\end{tikzpicture}
\end{center}}

\newcommand{\titrepage}{\begin{tcolorbox}[colback=white,colframe=bleu!70!black,boxrule=0.65pt,arc=3mm,left=2.5mm,right=2.5mm,top=1.15mm,bottom=1.15mm,before skip=0mm,after skip=1.4mm,drop fuzzy shadow=black!10]
{\Large\bfseries Parcours de révision -- Variables aléatoires}\hfill {\normalsize Première spécialité -- sans calculatrice}
\end{tcolorbox}}

\newcommand{\ficheRecap}{%
\begin{landscape}
\pagecolor{fondpage}
\thispagestyle{empty}
\titrepage
\begin{tabularx}{\linewidth}{@{}X@{\hspace{1.8mm}}X@{\hspace{1.8mm}}X@{}}
\begin{panel}{}
\tagcol{bleu}{Modéliser}
\begin{fichebox}{bleu}{Variable aléatoire}
Une variable aléatoire $X$ associe un nombre réel à chaque issue d'une expérience aléatoire.\\
Exemple : $X$ peut représenter un gain, un prix payé, un nombre de points, etc.
\end{fichebox}
\begin{fichebox}{bleu}{Valeurs prises}
On commence par lister toutes les valeurs possibles de $X$ :
\[
X(\Omega)=\{x_1;x_2;\ldots;x_k\}.
\]
Il ne faut pas oublier les valeurs négatives quand $X$ est un gain algébrique.
\end{fichebox}
\begin{fichebox}{bleu}{Loi de probabilité}
La loi de $X$ donne toutes les probabilités $P(X=x_i)$.
\loitableau
On doit toujours vérifier :
\[
p_1+p_2+\cdots+p_k=1.
\]
\end{fichebox}
\begin{fichebox}{bleu}{Gain algébrique}
Dans un jeu :
\[
\text{gain algébrique}=\text{somme reçue}-\text{prix payé}.
\]
C'est ce gain que l'on utilise pour décider si le jeu est favorable.
\end{fichebox}
\end{panel}
&
\begin{panel}{}
\tagcol{vert}{Espérance}
\begin{fichebox}{vert}{Formule}
Si $X$ prend les valeurs $x_1,\ldots,x_k$ avec les probabilités $p_1,\ldots,p_k$, alors :
\[
E(X)=p_1x_1+p_2x_2+\cdots+p_kx_k.
\]
\end{fichebox}
\begin{fichebox}{vert}{Interprétation}
$E(X)$ est la valeur moyenne attendue si l'on répète l'expérience un grand nombre de fois.
\end{fichebox}
\begin{fichebox}{vert}{Jeu équitable}
On note $X$ le \textbf{gain algébrique} du joueur.\\[0.3em]
\begin{center}
\renewcommand{\arraystretch}{1.15}
\begin{tabular}{c|c}
$E(X)>0$ & favorable au joueur\\
$E(X)=0$ & équitable\\
$E(X)<0$ & défavorable au joueur
\end{tabular}
\end{center}
\end{fichebox}
\begin{fichebox}{vert}{Phrase type}
\textit{Sur un grand nombre de parties, le gain moyen est d'environ $E(X)$ euro par partie.}
\end{fichebox}
\end{panel}
&
\begin{panel}{}
\tagcol{orange}{Variance et écart type}
\begin{fichebox}{orange}{Variance}
La variance mesure la dispersion autour de l'espérance.\\
Si $m=E(X)$, alors :
\[
V(X)=p_1(x_1-m)^2+p_2(x_2-m)^2+\cdots+p_k(x_k-m)^2.
\]
\end{fichebox}
\begin{fichebox}{orange}{Écart type}
\[
\sigma(X)=\sqrt{V(X)}.
\]
Plus l'écart type est grand, plus les valeurs de $X$ sont dispersées.
\end{fichebox}
\begin{fichebox}{orange}{Comparer deux jeux}
\begin{itemize}[leftmargin=*,itemsep=0.15em]
\item L'espérance compare le gain moyen.
\item La variance ou l'écart type compare le risque.
\item À espérance égale, le jeu de plus grand écart type est le plus risqué.
\end{itemize}
\end{fichebox}
\end{panel}
\end{tabularx}
\clearpage
\nopagecolor
\end{landscape}}

\begin{document}
\begin{center}
{\LARGE\bfseries Parcours de révision -- Variables aléatoires}\\[1mm]
{\large Première générale -- spécialité mathématiques}\\[1mm]
{\normalsize Sans calculatrice}
\end{center}
\tableofcontents\newpage
\phantomsection\addcontentsline{toc}{section}{Fiche récapitulative}
\ficheRecap

\section{Questions flash}
\setcounter{question}{0}
\begin{blocquestions}{Questions flash 1 à 12}
\item Une variable aléatoire $X$ prend les valeurs $0$, $1$ et $4$ avec les probabilités $\frac12$, $\frac14$ et $\frac14$. Vérifier que cela définit une loi de probabilité.
\item Pour la loi précédente, calculer $P(X=4)$.
\item Pour la loi précédente, calculer $P(X\geq 1)$.
\item Pour la loi précédente, calculer $E(X)$.
\item Une variable aléatoire $Y$ prend les valeurs $-2$ et $3$ avec probabilités $\frac35$ et $\frac25$. Calculer $E(Y)$.
\item Une variable aléatoire $X$ prend les valeurs $2$ et $8$ avec les probabilités $\frac14$ et $\frac34$. Calculer $E(X)$.
\item Une variable aléatoire $X$ prend les valeurs $0$ et $4$ avec les probabilités $\frac12$ et $\frac12$. Écrire la somme permettant de calculer $V(X)$.
\item Si $V(X)=9$, calculer $\sigma(X)$.
\item Dans un jeu, $X$ désigne le gain algébrique et $E(X)=0$. Que peut-on dire du jeu ?
\item Dans un jeu, $E(X)=-0,5$. Le jeu est-il favorable au joueur ?
\item Donner la formule de l'espérance d'une variable aléatoire prenant les valeurs $x_i$ avec probabilités $p_i$.
\item Donner la formule explicite de la variance à partir de l'espérance $m=E(X)$.
\end{blocquestions}
\begin{blocquestions}{Questions flash 13 à 24}
\item On lance une pièce équilibrée. $X=1$ si on obtient pile et $X=0$ sinon. Donner la loi de $X$.
\item Pour la variable précédente, calculer $E(X)$.
\item On lance un dé équilibré. $X$ vaut $1$ si on obtient un nombre pair, et $0$ sinon. Donner $P(X=1)$.
\item Une urne contient 2 boules rouges et 3 boules bleues. On gagne 4 euros si on tire rouge, sinon on perd 1 euro. Donner les valeurs prises par le gain $G$.
\item Pour la situation précédente, donner la loi de $G$.
\item Pour la situation précédente, calculer $E(G)$.
\item Un jeu coûte 3 euros. On reçoit 10 euros avec probabilité $\frac15$ et rien sinon. Donner les valeurs prises par le gain algébrique $X$.
\item Pour la situation précédente, calculer $E(X)$.
\item Une loi donne $P(X=0)=\frac14$, $P(X=2)=\frac14$ et $P(X=5)=\frac12$. Calculer $E(X)$.
\item Avec la loi précédente, écrire la somme permettant de calculer $V(X)$.
\item Avec la loi précédente, calculer $V(X)$.
\item Si deux jeux ont la même espérance mais pas le même écart type, lequel est le plus risqué ?
\end{blocquestions}
\begin{blocquestions}{Questions flash 25 à 36}
\item Une variable $X$ prend les valeurs $1$ et $5$ avec probabilités $\frac12$ et $\frac12$. Calculer $E(X)$.
\item Pour la variable précédente, calculer $V(X)$.
\item Une remise vaut 0 euro avec probabilité $\frac12$ et 10 euros avec probabilité $\frac12$. Calculer la remise moyenne.
\item Un prix initial vaut 40 euros. La remise moyenne vaut 6 euros. Quel est le prix moyen payé ?
\item Si $X$ est un gain et $Y=X-2$, que représente souvent le nombre $2$ dans un jeu ?
\item Une variable $X$ prend les valeurs $2$ et $6$ avec probabilités $\frac12$ et $\frac12$. Calculer $E(X)$.
\item Si $\sigma(X)=2$, que vaut $V(X)$ ?
\item Si $V(X)=0$, que peut-on dire de $X$ ?
\item Une variable $X$ prend les valeurs $-1$, $0$, $2$ avec probabilités $\frac14$, $\frac12$, $\frac14$. Calculer $E(X)$.
\item Pour cette même variable, écrire la somme permettant de calculer $V(X)$.
\item Pour cette même variable, calculer $V(X)$.
\item En une phrase, expliquer ce que mesure l'écart type.
\end{blocquestions}

\section{QCM d'automatismes}
Pour chaque question, une seule réponse est correcte.
\begin{enumerate}
\item Une loi de probabilité doit vérifier :
\qcm{la somme des probabilités vaut $0$}{la somme des probabilités vaut $1$}{toutes les valeurs sont positives}{l'espérance vaut $1$}
\item Si $X$ prend les valeurs $0$ et $10$ avec probabilités $\frac12$ et $\frac12$, alors $E(X)=$
\qcm{$0$}{$5$}{$10$}{$20$}
\item Si $X$ prend les valeurs $1$ et $3$ avec probabilités $\frac12$ et $\frac12$, alors $E(X)=$
\qcm{$1$}{$2$}{$3$}{$4$}
\item Si $X$ prend les valeurs $0$ et $4$ avec probabilités $\frac12$ et $\frac12$, alors $V(X)=$
\qcm{$2$}{$4$}{$8$}{$16$}
\item L'écart type est égal à :
\qcm{$V(X)^2$}{$\sqrt{V(X)}$}{$E(X)^2$}{$E(X^2)$}
\item Un jeu équitable vérifie :
\qcm{$E(X)=0$}{$V(X)=0$}{$E(X)=1$}{$\sigma(X)=0$}
\item Si $E(X)<0$ pour un gain algébrique, alors le jeu est :
\qcm{favorable au joueur}{défavorable au joueur}{forcément équitable}{impossible}
\item La variance mesure surtout :
\qcm{une probabilité}{une fréquence}{une dispersion}{une valeur minimale}
\end{enumerate}

\section{Sujets type bac / E3C}

\titreSujet{Sujet 1 -- Jeu de jetons et espérance}{inspiré d'exercices E3C sur variable aléatoire, gain algébrique et jeu équitable}
Une urne contient 10 jetons indiscernables au toucher : 1 jeton rouge, 3 jetons bleus et 6 jetons blancs.
Un joueur tire un jeton au hasard. Le jeu coûte 2 euros.
\begin{itemize}
\item S'il tire le jeton rouge, il reçoit 8 euros.
\item S'il tire un jeton bleu, il reçoit 3 euros.
\item S'il tire un jeton blanc, il ne reçoit rien.
\end{itemize}
On note $X$ le gain algébrique du joueur, c'est-à-dire la somme reçue moins le prix du jeu.
\begin{enumerate}
\item Déterminer les valeurs prises par $X$.
\item Déterminer la loi de probabilité de $X$ sous forme de tableau.
\item Calculer l'espérance de $X$.
\item Interpréter ce résultat dans le contexte.
\item Le jeu est-il favorable au joueur ?
\item Quel prix d'entrée faudrait-il choisir pour que le jeu soit équitable ?
\end{enumerate}

\titreSujet{Sujet 2 -- Deux jeux à comparer}{inspiré d'exercices E3C sur espérance, variance et comparaison du risque}
On propose deux jeux. Dans chaque jeu, on lance un dé équilibré à six faces.

\textbf{Jeu A :}
si le résultat est 6, on gagne 9 euros ; si le résultat est 4 ou 5, on gagne 3 euros ; sinon, on perd 3 euros.

\textbf{Jeu B :}
si le résultat est pair, on gagne 3 euros ; si le résultat est impair, on perd 3 euros.

On note $A$ la variable aléatoire donnant le gain au jeu A et $B$ la variable aléatoire donnant le gain au jeu B.
\begin{enumerate}
\item Déterminer la loi de probabilité de $A$.
\item Déterminer la loi de probabilité de $B$.
\item Calculer $E(A)$ et $E(B)$.
\item Calculer $V(A)$ et $V(B)$.
\item En déduire les écarts types de $A$ et $B$.
\item Les deux jeux ont-ils le même gain moyen ?
\item Quel jeu paraît le plus risqué ? Justifier.
\end{enumerate}

\titreSujet{Sujet 3 -- Remise commerciale et prix payé}{inspiré d'exercices E3C de modélisation par variable aléatoire en contexte économique}
Un magasin organise une opération promotionnelle. À chaque passage en caisse, le client tire une carte au hasard parmi 20 cartes indiscernables.
\begin{itemize}
\item 2 cartes donnent une remise de 20 euros ;
\item 5 cartes donnent une remise de 10 euros ;
\item les autres cartes ne donnent aucune remise.
\end{itemize}
Un client achète un article coûtant 50 euros. On note $P$ la variable aléatoire égale au prix réellement payé par le client après remise.
\begin{enumerate}
\item Déterminer les valeurs prises par $P$.
\item Déterminer la loi de probabilité de $P$.
\item Calculer $E(P)$.
\item Interpréter $E(P)$ dans le contexte.
\item Calculer $V(P)$.
\item Calculer l'écart type de $P$.
\item Le magasin affirme : « en moyenne, chaque client bénéficie de 5 euros de remise ». Cette affirmation est-elle vraie ?
\end{enumerate}

\end{document}
